% Lösungen Einheit 3 von 3 – Kreisteile (zusammenbau v0.1)
\einheitenkopf{Einheit 3 von 3 · Kreisteile}

%% TODO Lösung der Erklärzeile 20g) fehlt in der Bank
\erg{20}{a) $\frac{1}{2}$ \quad b) $\frac{1}{4}$ \quad c) $\frac{1}{2}$ \quad d) $\frac{1}{6}$ \quad e) $\frac{1}{6}$ \quad f) $\frac{1}{4}$ \quad g) \ldots \quad h) $100 : 360 \approx 0{,}2778 \approx 27{,}8\,\%$ \quad i) $\alpha = \frac{3}{8} \cdot 360^\circ = 135^\circ$ \quad j) $b = \frac{40}{360} \cdot 2 \cdot \pi \cdot 9 \approx 6{,}3$ cm \quad k) $A = \frac{84}{360} \cdot \pi \cdot 6^2 \approx 26{,}4$ cm² \quad l) $b = \frac{45}{360} \cdot 2 \cdot \pi \cdot 8 \approx 6{,}28$ cm; Umfang $\approx 6{,}28 + 2 \cdot 8 \approx 22{,}3$ cm \quad m) $A = \pi \cdot 7^2 - \pi \cdot 4^2 \approx 103{,}7$ cm² \quad n) $130 : 360 \approx 0{,}3611 \approx 36{,}1\,\%$}
\erg{21}{a) $A = \frac{150}{360} \cdot \pi \cdot 6^2 \approx 47{,}1$ m²}
\erg{22}{a) Jana hat mit dem Winkel des weißen Teils gerechnet. Richtig: $110 : 360 \approx 0{,}306 \approx 30{,}6\,\%$ \quad b) Die beiden Radien des Halbkreises liegen auf einer Geraden, dem Durchmesser; das ist ein gestreckter Winkel, also die Hälfte des Vollwinkels. \quad c) $A = \frac{1}{4} \cdot \pi \cdot 7^2 \approx 38{,}48$ m²; $38{,}48 \cdot 4 \approx 153{,}9$ Liter \quad d) Bus $= 45\,\%$, Rad $= 30\,\%$, Auto $= 25\,\%$}
